% Complete Section 7.4, original PDF 96-98 / printed 86-88.
% block: finite-element-differential-forms-heading page:96 type:section
\section{{有限元微分形式}}\label{sec:finite-element-differential-forms}

% block: complete-form-degrees-of-freedom page:96 type:ordinary-paragraph
{现在可以定义我们的有限元空间. 形函数取自某个完全多项式空间 $\mathcal P_r\Lambda^k(T)$, 或某个修剪多项式空间 $\mathcal P_r^-\Lambda^k(T)$. 一个重要观察是, 两类空间彼此相关, 放在一起处理最易理解. 事实上, 完全空间的自由度涉及面上的修剪空间, 反之亦然. 具体地, $\omega\in\mathcal P_r\Lambda^k(T)$ 的自由度为}
\begin{equation}\label{eq:complete-form-degrees-of-freedom}
\omega\mapsto\int_f(\operatorname{tr}_f\omega)\wedge\mu,\qquad\mu\in\mathcal P_{r+k-d}^-\Lambda^{d-k}(f),\quad f\in\Delta_d(T),\ d\geq k,
\end{equation}

% block: trimmed-form-degrees-of-freedom page:97 type:ordinary-paragraph
{而 $\omega\in\mathcal P_r^-\Lambda^k(T)$ 的自由度为}
\begin{equation}\label{eq:trimmed-form-degrees-of-freedom}
\omega\mapsto\int_f(\operatorname{tr}_f\omega)\wedge\mu,\qquad\mu\in\mathcal P_{r+k-d-1}\Lambda^{d-k}(f),\quad f\in\Delta_d(T),\ d\geq k.
\end{equation}
{注意, 每个自由度都与单纯形 $T$ 的某个特定面 $f$ 关联, 其维数介于 $0$ 与 $n$ 之间 ($\Delta_d(T)$ 表示 $T$ 的所有 $d$ 维面的集合). 然后要求自由度单值, 就确定了组装后的有限元空间.}

% block: form-unisolvence-goal page:97 type:ordinary-paragraph
{为使自由度规定 \eqref{eq:complete-form-degrees-of-freedom} 和 \eqref{eq:trimmed-form-degrees-of-freedom} 良好地定义有限元空间, 必须分别验证其\emph{唯一可解性}. 这需要证明: 首先, 关联于形函数空间的自由度总数不超过空间维数; 其次, 若形函数空间中某个 $\omega$ 的所有自由度都为零, 则 $\omega$ 本身为零. 对两种形函数空间都可以做到这一点. 这里对任意多项式次数 $r\geq1$, 任意维数 $n\geq1$ 和任意形式次数 $0\leq k\leq n$, 验证修剪空间 $\mathcal P_r^-\Lambda^k(T)$ 的唯一可解性. 证明基于 Koszul 微分的性质, 直接取自 \BookCite{6}. 完全多项式空间的情形有十分相似的证明.}

% block: form-degree-of-freedom-count page:97 type:ordinary-paragraph
{\eqref{eq:trimmed-form-degrees-of-freedom} 定义的自由度数为}
\[
\begin{aligned}
&\sum_{d\geq k}\#\Delta_d(T)\dim[\mathcal P_{r+k-d-1}\Lambda^{d-k}(\mathbb R^d)]\\
&\qquad=\sum_{d\geq k}\binom{n+1}{d+1}\binom{r+k-1}d\binom d{d-k}\\
&\qquad=\sum_{j\geq0}\binom{n+1}{j+k+1}\binom{r+k-1}{j+k}\binom{j+k}j.
\end{aligned}
\]
{利用二项式恒等式化简,}
\[
\binom ab\binom bc=\binom ac\binom{a-c}{a-b}
\]
{(它直接来自二项式系数的定义), 以及 Vandermonde 恒等式}
\[
\sum_{j\geq0}\binom a{b+j}\binom cj=\binom{a+c}{a-b},
\]
{右端变为}
\[
\binom{r+n}{r+k}\binom{r+k-1}k=\dim\mathcal P_r^-\Lambda^k(T),
\]
{如所要求.}

% block: form-unisolvence-induction pages:97,98 type:ordinary-paragraph
{还需证明, 若某个 $\omega\in\mathcal P_r^-\Lambda^k(T)$ 的所有自由度都为零, 则 $\omega$ 为零. 对维数 $n$ 作归纳证明 (情形 $n=1$ 容易). 对 $T$ 的任意超面 (余维为 $1$ 的面) $f$, $\operatorname{tr}_f\omega$ 属于 $\mathcal P_r^-\Lambda^k(f)$, 且 $\operatorname{tr}_f\omega$ 的自由度为零, 因为它们是 $\omega$ 的自由度的一部分. 因此可用归纳假设得出, 对所有超面 $f$, $\operatorname{tr}_f\omega$ 都为零, 即 $\omega\in\mathring{\mathcal P}_r^-\Lambda^k(T)$, 这是 $\mathcal P_r^-\Lambda^k(T)$ 中在整个边界上迹为零的元素构成的子空间. 由于迹算子与外微分交换 (任意拉回都如此), 有 $\nu:=d\omega\in\mathring{\mathcal P}_{r-1}\Lambda^{k+1}(T)$. 此外,}
\[
\int_T\nu\wedge p=\pm\int_T\omega\wedge dp=0,\qquad p\in\mathcal P_{r+k-n}\Lambda^{n-k-1}(T),
\]
{其中第一个等式来自分部积分, 第二个来自 $\omega$ 的自由度为零. 由此可以通过基于重心坐标的具体构造 [\BookNumber{6}, Lemma 3.4], 证明 $\nu=0$. 由于 $\omega\in\mathcal P_r^-\Lambda^k$ 且 $d\omega=0$, 可得 $\omega\in\mathcal P_{r-1}\Lambda^k$, 再结合其迹为零, 意味着 $\omega\in\mathring{\mathcal P}_{r-1}\Lambda^k$. 再次应用 [\BookNumber{6}, Lemma 3.4] 的构造, 表明 $\omega$ 为零, 从而完成唯一可解性的证明.}

% block: assembled-form-families page:98 type:ordinary-paragraph
{由此在每个单元 $T$ 上定义了一族用作形函数的多项式微分 $k$-形式空间, 即对每个 $r\geq1$ 定义的空间 $\mathcal P_r\Lambda^k(T)$ 和 $\mathcal P_r^-\Lambda^k(T)$. 此外, 对每个空间都指定了一组具有唯一可解性的自由度. 结合剖分 $\mathcal T_h$, 每一种这样的选择都确定相应的分片多项式微分形式有限元空间 $\mathcal P_r\Lambda^k(\mathcal T_h)$ 或 $\mathcal P_r^-\Lambda^k(\mathcal T_h)$.}

% block: assembled-form-continuity-projections page:98 type:ordinary-paragraph
{\eqref{eq:complete-form-degrees-of-freedom} 和 \eqref{eq:trimmed-form-degrees-of-freedom} 定义的自由度具有许多良好性质, 它们反映在组装后有限元空间的性质中. 特别地,}
\begin{itemize}
\item{这些自由度恰好施加了组装后的有限元函数属于 $H\Lambda^k$ 所需的连续性,}
\[
\mathcal P_r\Lambda^k(\mathcal T_h)=\{\,\omega\in H\Lambda^k(\Omega)\mid\omega|_T\in\mathcal P_r\Lambda^k(T)\,\};
\]
\item{自由度确定的典范投影算子 $\Pi_{r,h}^k:C\Lambda^k(\bar\Omega)\to\mathcal P_r\Lambda^k(\mathcal T_h)$ 构成上链映射; 对 $\mathcal P^-$ 空间及其典范投影 $\Pi_{r,h}^{-,k}$ 也有类似结论.}
\end{itemize}

% block: canonical-projection-commutativity page:98 type:ordinary-paragraph
{典范投影与外微分的交换性, 可用 Stokes 定理验证. 这是一个关键性质: 它是构造有界上链投影的主要步骤.}
